\section{gpt pro}
\subsection*{When does $\dA=\po$ almost surely?}

Recall the definition
\[
    \dA \;\triangleq\; {F^\star}'\!\left(F'(a)-\eta\,\hl\right),
    \qquad
    \po \;\triangleq\; \E\!\left[A_\star \,\middle|\, A, Y\right],
\]
and (in the multiplicative-$Y$ class) $\hl=\frac{Y}{\eta P_A}T_A$ with $T\in(\R^d)^{\cA}$.

\begin{corollary}[Quadratic potential $\Rightarrow$ linear regression in $Y$]
    \label{cor:quad-linear}
    Let $F(x)=\tfrac12\|x\|^2$. Then $\dA=\po$ a.s.\ iff, for $\PP$-a.e.\ $b$,
    \begin{equation}
        \label{eq:lin-reg}
        \E[A_\star\mid A=b,Y=y] \;=\; a - \gamma_b\,y
        \quad\text{for $Y$-a.e.\ $y$,}
    \end{equation}
    with some $\gamma_b\in\R^{d}$, in which case choosing $T_b=P_b\,\gamma_b$ yields $\dA=\po$ a.s.
    A canonical sufficient condition for \eqref{eq:lin-reg} is joint conditional Gaussianity:
    if $(A_\star,Y)\mid(A=b)$ is jointly Gaussian with
    \(
    \E[A_\star\mid A=b]=\pr,\;
    \E[Y\mid A=b]=\mu_b,
    \)
    and
    \(
    \Gamma_b=\Cov(A_\star,Y\mid A=b),\;
    \sigma_b^2=\Var(Y\mid A=b)>0,
    \)
    then
    \[
        \po \;=\; \pr + \Gamma_b\,\sigma_b^{-2}\,(Y-\mu_b).
    \]
    Setting
    \(
    T_b \;=\; -P_b\,\Gamma_b\,\sigma_b^{-2}
    \)
    and
    \(
    a \;=\; \pr - \Gamma_b\,\sigma_b^{-2}\,\mu_b
    \)
    forces \eqref{eq:lin-reg} and hence $\dA=\po$ a.s.
    In particular, if $\E[Y\mid A=b]=0$ (e.g.\ $\E[\ell]=0$ under $\nu$), it suffices to take $a=\pr$ and $T_b=-P_b\,\Gamma_b\,\sigma_b^{-2}$.
\end{corollary}

\begin{corollary}[Mirror–exponential-family conjugacy]
    \label{cor:mirror-conjugacy}
    Assume $F$ is Legendre and consider the exponential family on $\cA$ with log-partition $F^\star$:
    \(
    \pi_\theta(\d x)\propto h(x)\,\exp\{\langle\theta,x\rangle - F^\star(\theta)\}\,\d x.
    \)
    Suppose $A_\star\sim \pi_{\theta_0}$ with $\theta_0=F'(a)$, and, for each $b$,
    the conditional law of $Y$ given $(A_\star=x,A=b)$ has density
    \[
        q_b(y\mid x) \;\propto\; r_b(y)\,\exp\!\left\{-\tfrac{y}{P_b}\,\langle T_b, x\rangle\right\},
    \]
    for some measurable $r_b$ and fixed $T_b\in\R^d$.
    Then Bayes' rule gives the posterior
    \(
    A_\star\mid (A=b,Y=y)\sim \pi_{\theta_0-(y/P_b)T_b},
    \)
    hence
    \(
    \po={F^\star}'\!\big(\theta_0-(y/P_b)T_b\big)
    ={F^\star}'\!\left(F'(a)-\tfrac{y}{P_b}T_b\right)=\dA
    \)
    a.s.
\end{corollary}

\begin{corollary}[Degenerate $Y$]
    \label{cor:degenerate}
    If $Y$ is (conditionally) degenerate given $A=b$ (i.e.\ $\PP(Y=y_b\mid A=b)=1$), then there exists $T_b$ with
    \(
    F'\!\big(\po\big)=F'(a)-\tfrac{y_b}{P_b}T_b
    \)
    and consequently $\dA=\po$ a.s.\ (this is vacuous for learning, but it shows attainability in degenerate cases).
\end{corollary}

\paragraph{A necessary condition (no hidden curvature in $y$).}
\begin{lemma}[Affine-dual necessity]
    \label{lem:affine-necessary}
    Suppose $E_T$ is multiplicative-$Y$ and $\dA=\po$ a.s.
    Fix $b$ with $\PP(A=b)>0$ and assume $y\mapsto \po(b,y)$ is a.s.\ differentiable.
    Then the \emph{dual} posterior mean must be affine in $y$:
    \[
        \frac{\partial^2}{\partial y^2}\,F'\!\big(\po(b,y)\big) \;=\; 0
        \quad\text{for $Y$-a.e.\ $y$.}
    \]
    For $F(x)=\tfrac12\|x\|^2$, this reduces to the linear-regression property
    $\frac{\partial^2}{\partial y^2}\E[A_\star\mid A=b,Y=y]=0$ a.s.
\end{lemma}

\paragraph{Why the condition is restrictive.}
In general, the map $y\mapsto \po$ need \emph{not} be affine (in primal or dual coordinates).%
\footnote{See the counterexample constructed at the end of the document, where the posterior mean contains a quadratic term in the latent $Z$; there the ratio $\E[D_F(\dA,\pr)]/\E[D_F(\po,\pr)]$ can be driven to $0$, showing that even a uniform multiplicative bound fails in general.}
This explains why the equality $\dA=\po$ cannot be enforced by multiplicative-$Y$ estimators unless one of the structural scenarios above holds.

\paragraph{Summary.}
\begin{itemize}
    \item With the \emph{unrestricted} optimal estimator $G_{\nu,P}$, one always has $\dA=\po$ a.s.
    \item Within the \emph{multiplicative-$Y$} class, $\dA=\po$ a.s.\ \emph{iff} the conditional posterior mean evolves along a straight line in \emph{dual} space as $y$ varies (Prop.~\ref{prop:multY-iff}). This is guaranteed by (i) quadratic $F$ $+$ linear conditional regression (e.g.\ conditional Gaussianity) or (ii) mirror–exponential-family conjugacy where observing $Y$ induces an additive shift of the natural parameter.
    \item Absent such structure, $\dA=\po$ a.s.\ typically fails.
\end{itemize}


\hrule

\subsubsection{Optimizing for $E$---linear (affine) in $Y$}
Motivated by the multiplicative class, we now allow an intercept (bias).  For
matrices $T,S \in \R^{|\cA| \times d}$ define the (affine) estimator
\[
    E_{T,S}(a,y)
    \;=\;
    \frac{y}{\eta P_a}\,T_a \;+\; S_a,
    \qquad
    \hl \;=\; E_{T,S}(A,Y),
\]
so $\eta \hl = \frac{Y}{P_A}T_A + \eta S_A$.  We still write
$\tilde{A} \triangleq {F^\star}'\!\big(F'(a)-\eta \hl\big)$ and use
$\Lambda(P,E\Vert \nu)$ from \eqref{eq:Lambda-simplified}.

\paragraph{First-order optimality with slope \& intercept.}
For fixed $a,\eta,P,\nu$, consider
\(
\inf_{T,S} \Lambda(P, E_{T,S}\Vert \nu)
\).
The next theorem gives the optimality conditions.

\begin{theorem}[Linear-$Y$ optimality conditions]
    \label{thm:optimal-linY}
    Let $(T^\star,S^\star)$ minimize $\Lambda(P,E_{T,S}\Vert \nu)$ over the affine
    class and set $\hl=E_{T^\star,S^\star}(A,Y)$ and
    $\tilde{A}={F^\star}'(F'(a)-\eta \hl)$. Then for every $b\in\cA$,
    \begin{align}
        \E\!\left[\, Y\,(A_\star-\tilde{A}) \,\big|\, A=b \right] & = 0,
        \label{eq:linY-FOC-slope}
        \\
        \E\!\left[\, A_\star-\tilde{A} \,\big|\, A=b \right]      & = 0.
        \label{eq:linY-FOC-intercept}
    \end{align}
    Equivalently, conditional on $A=b$, the residual $A_\star-\tilde{A}$ is
    orthogonal to both the constant function $1$ and to the scalar $Y$.
\end{theorem}

\begin{proof}
    Starting from \eqref{eq:Lambda-simplified} and conditioning on $A=b$,
    minimize
    \[
        \Phi_b(T_b,S_b)
        =
        \E\!\left[
            \left\langle A_\star, \tfrac{Y}{\eta P_b}T_b + S_b \right\rangle
            + \tfrac{1}{\eta} F^\star\!\left(F'(a)-\tfrac{Y}{P_b}T_b-\eta S_b\right)
            \,\Big|\, A=b
            \right].
    \]
    Differentiating and setting the gradients to zero yields
    \[
        0=\tfrac{\partial \Phi_b}{\partial T_b}
        \;\Rightarrow\;
        \E\!\left[ Y A_\star - Y\,{F^\star}'\!\left(F'(a)-\tfrac{Y}{P_b}T_b-\eta S_b\right)\Big|A=b\right]=0,
    \]
    and
    \[
        0=\tfrac{\partial \Phi_b}{\partial S_b}
        \;\Rightarrow\;
        \E\!\left[ A_\star - {F^\star}'\!\left(F'(a)-\tfrac{Y}{P_b}T_b-\eta S_b\right)\Big|A=b\right]=0.
    \]
    With $\hl=\frac{Y}{\eta P_b}T_b+S_b$ these are exactly
    \eqref{eq:linY-FOC-slope}--\eqref{eq:linY-FOC-intercept}.
\end{proof}

\paragraph{Two immediate consequences.}
\begin{lemma}[Calibration and replacement identity]
    \label{lem:linY-calibration-replace}
    Under the conditions of Theorem~\ref{thm:optimal-linY}:
    \begin{enumerate}
        \item (\emph{Calibration}) Since $A$ is independent of $(A_\star,\ell)$ under $\nu$,
              \eqref{eq:linY-FOC-intercept} implies
              \[
                  \E[\tilde{A}\mid A=b] \;=\; \E[A_\star\mid A=b] \;=\; \pr,\qquad\forall b\in\cA.
              \]
        \item (\emph{Replacing $A_\star$ by $\tilde{A}$ inside $\langle\cdot,\hl\rangle$})
              Writing $\hl=\frac{Y}{\eta P_A}T_A + S_A$ and taking conditional expectations,
              \eqref{eq:linY-FOC-slope}--\eqref{eq:linY-FOC-intercept} give
              \[
                  \E\!\left[\langle A_\star,\hl\rangle\mid A\right]
                  \;=\;
                  \E\!\left[\langle \tilde{A},\hl\rangle\mid A\right].
              \]
    \end{enumerate}
\end{lemma}

\begin{lemma}[E-optimized $\Lambda$ (linear-$Y$)]
    \label{lem:linY-optimized-Lambda}
    Let $(T^\star,S^\star)$ be optimal and $\hl=E_{T^\star,S^\star}(A,Y)$. Then
    \begin{equation}
        \Lambda(P,E_{T^\star,S^\star}\Vert \nu)
        \;=\;
        \frac{1}{\eta}\,\Delta(\nu,P)
        \;-\;
        \frac{1}{\eta^2}\,
        \E\!\left[
            D_{F^\star}\!\left(F'(a),\,F'(a)-\eta \hl\right)
            \right].
        \label{eq:linY-optimized-Lambda}
    \end{equation}
\end{lemma}

\begin{proof}
    Start from \eqref{eq:Lambda-simplified}. By
    Lemma~\ref{lem:linY-calibration-replace}(ii) we can replace
    $\langle A_\star,\hl\rangle$ by $\langle \tilde{A},\hl\rangle$ inside the
    expectation. The same algebra as in Lemma~\ref{lem:E-optimized-Lambda} then
    yields \eqref{eq:linY-optimized-Lambda}.
\end{proof}

\paragraph{Orthogonality and an expected-Bregman identity (unchanged).}
Using $F'(a)-F'(\tilde{A})=\eta \hl$ and
Lemma~\ref{lem:linY-calibration-replace}, we get
\[
    \E\!\left[
        \big\langle F'(a)-F'(\tilde{A}),\, \po-\tilde{A}\big\rangle
        \,\Big|\, A
        \right]
    = \eta\,\E\!\left[\langle \hl,\, \po-\tilde{A}\rangle \big| A\right]
    = 0.
\]
Applying the three-point identity for Bregman divergences and averaging over $A$
gives exactly the same decomposition as in
\eqref{eq:expected-bregman-identity}:
\begin{equation}
    \E\!\left[D_F(\po,\tilde{A}) + D_F(\tilde{A},a)\right]
    \;=\;
    \E\!\left[D_F(\po,\pr)\right] + D_F(\pr,a).
    \label{eq:expected-bregman-identity-linY}
\end{equation}
In particular, conditioning on $A$ and using the calibration
$\E[\tilde{A}\mid A]=\pr$ yields
\begin{equation}
    \E\!\left[D_F(\tilde{A},a)\right]
    \;=\;
    \E\!\left[D_F(\tilde{A},\pr)\right] + D_F(\pr,a),
    \qquad
    \E\!\left[D_F(\po,\tilde{A})\right] + \E\!\left[D_F(\tilde{A},\pr)\right]
    \;=\;
    \E\!\left[D_F(\po,\pr)\right].
    \label{eq:linY-two-way}
\end{equation}

\paragraph{Implications for information-ratio style bounds.}
The optimized value \eqref{eq:linY-optimized-Lambda} has the \emph{same form}
as in the multiplicative-only class. To relate it to the information ratio as in
Theorem~\ref{thm:optimal-Lambda-info-ratio}, one would want
$\E[D_F(\tilde{A},a)]$ to be comparable to $\E[D_F(\po,\pr)]$.  From
\eqref{eq:expected-bregman-identity-linY}--\eqref{eq:linY-two-way},
\[
    \E[D_F(\tilde{A},a)]
    \;=\;
    \E[D_F(\po,\pr)] - \E[D_F(\po,\tilde{A})] + D_F(\pr,a),
\]
so any universal lower bound of the form
\(
\E[D_F(\tilde{A},a)] \ge \alpha\,\E[D_F(\po,\pr)]
\)
would require $\E[D_F(\po,\tilde{A})]\le (1-\alpha)\,\E[D_F(\po,\pr)]$.
This fails in general, even with the bias term, as the next corollary notes.

\begin{corollary}[No universal multiplicative factor, even with bias]
    \label{cor:no-alpha-linY}
    The negative example given under \emph{Multiplicative factor?} continues to
    apply to the linear-$Y$ class.  In that construction with quadratic $F$,
    $\E[Y\mid A]=0$, so the intercept condition \eqref{eq:linY-FOC-intercept}
    imposes no constraint on $S$, while the slope condition
    \eqref{eq:linY-FOC-slope} uniquely determines $T$ exactly as in the
    multiplicative case (and $S$ cancels from \eqref{eq:linY-FOC-slope}).  Hence
    $\E[D_F(\tilde{A},\pr)]/\E[D_F(\po,\pr)]$ can still be driven arbitrarily close
    to $0$ by taking $\omega\to\infty$, and there is no universal
    $\alpha>0$ valid for all instances.
\end{corollary}

\paragraph{When can we achieve $\tilde{A}=\po$ a.s.?}
Allowing a bias relaxes the intercept but does not remove the fundamental
nonlinearity barrier.

\begin{proposition}[Characterization in the linear-$Y$ class]
    \label{prop:linY-iff}
    Consider $E_{T,S}$ with $\hl=\frac{Y}{\eta P_A}T_A + S_A$ and
    $\tilde{A}={F^\star}'(F'(a)-\eta \hl)$. Then $\tilde{A}=\po$ a.s.\ if and only if
    \begin{equation}
        \label{eq:affine-dual-linY}
        F'\!\big(\po\big)
        \;=\;
        F'(a) \;-\; \frac{Y}{P_A}\,T_A \;-\; \eta\,S_A
        \qquad\text{almost surely.}
    \end{equation}
    Equivalently, conditional on $A=b$, the map
    $y\mapsto F'\!\big(\E[A_\star\mid A=b,Y=y]\big)$ must be \emph{affine in $y$},
    with slope $-T_b/P_b$ and intercept $F'(a)-\eta S_b$.
\end{proposition}

\begin{proof}
    By definition $\tilde{A}=\po$ a.s.\ $\iff$
    \(
    \po={F^\star}'\!\big(F'(a)-\eta \hl\big)
    \)
    a.s.\ Applying $F'$ (injective for Legendre $F$) and substituting
    $\eta \hl=\frac{Y}{P_A}T_A+\eta S_A$ gives \eqref{eq:affine-dual-linY}.
\end{proof}

\begin{remark}[What changed with the bias?]
    Compared to the multiplicative class (Prop.~\ref{prop:mult-Y-iff}), the bias
    frees the \emph{intercept} in \eqref{eq:affine-dual-linY} (from $F'(a)$ to
    $F'(a)-\eta S_A$) and enforces the calibration
    $\E[\tilde{A}\mid A]=\pr$.  However, unless $y\mapsto F'(\po)$ is already affine
    (in which case we can pick $T,S$ to match it), we still cannot force
    $\tilde{A}=\po$ a.s., and the value \eqref{eq:linY-optimized-Lambda} (hence the
    regret bound driven by it) remains of the same form as in the multiplicative
    case.
\end{remark}
